\documentclass[11pt,a4paper]{article}

\usepackage[margin=1.4cm,top=1.6cm,bottom=1.6cm]{geometry}
\usepackage[T1]{fontenc}
\usepackage[utf8]{inputenc}
\usepackage{amsmath,amssymb,amsfonts}
\usepackage{siunitx}
\usepackage{booktabs}
\usepackage{tabularx}
\usepackage{array}
\usepackage{enumitem}
\usepackage{fancyhdr}
\usepackage{xcolor}
\usepackage{microtype}
\linespread{0.97}
\usepackage{lastpage}

% Page styling and headers
\pagestyle{fancy}
\fancyhf{}
\fancyhead[L]{\small\textbf{NIT Calicut} \textbullet\ \textit{Department of Physics}}
\fancyhead[R]{\small\textbf{PH7302E: Condensed Matter Physics}}
\fancyfoot[L]{\small\textit{End-Semester Examination --- 2026--27}}
\fancyfoot[C]{Page \thepage\ of \pageref{LastPage}}
\fancyfoot[R]{\small\textit{M.Sc. Physics (2023 Curriculum)}}
\renewcommand{\headrulewidth}{0.4pt}
\renewcommand{\footrulewidth}{0.4pt}

\newcommand{\qmarks}[1]{\hfill\textbf{[#1~Marks]}}

\begin{document}

% -------------------------------------------------------------
% Header Block
% -------------------------------------------------------------
\noindent
\textbf{Course Code:} PH7302E \hfill \textbf{Session:} Winter Semester 2026--27 \\
\textbf{Course Title:} Condensed Matter Physics \hfill \textbf{Duration:} 3 Hours \\
\textbf{Max. Marks:} 50 \hfill \textbf{Weightage:} Pre-Midterm 30\% (15 M) / Post-Midterm 70\% (35 M)
%\vspace{2pt}
\rule{\textwidth}{0.5pt}

%\vspace{3pt}
\noindent\textbf{Instructions to Candidates:}
\begin{enumerate}[topsep=1pt,itemsep=0.5pt,parsep=0pt,leftmargin=1.4em]
    \item Answer \textbf{ALL} questions. Total marks: \textbf{50}.
    \item The examination is structured into two parts:
    \begin{itemize}[leftmargin=1.4em,itemsep=0pt]
        \item \textbf{PART A (Questions 1--2):} Pre-Midterm Syllabus \hfill \textbf{[15 Marks --- 30\%]}
        \item \textbf{PART B (Questions 3--6):} Post-Midterm Syllabus \hfill \textbf{[35 Marks --- 70\%]}
    \end{itemize}
    \item Marks for each question and its sub-parts are indicated in the right margin.
    \item State clearly any physical approximations and assumptions made.
 %   \item Non-programmable scientific calculators are permitted.
\end{enumerate}

\vspace{2pt}
\noindent\textbf{Physical Constants:}
$h = 6.626 \times 10^{-34}\text{ J}\cdot\text{s}$, $\hbar = 1.055 \times 10^{-34}\text{ J}\cdot\text{s}$, $e = 1.602 \times 10^{-19}\text{ C}$, $m_e = 9.109 \times 10^{-31}\text{ kg}$, $k_B = 1.381 \times 10^{-23}\text{ J}\cdot\text{K}^{-1}$, $N_A = 6.022 \times 10^{23}\text{ mol}^{-1}$, $\epsilon_0 = 8.854 \times 10^{-12}\text{ F}\cdot\text{m}^{-1}$, $\mu_0 = 4\pi \times 10^{-7}\text{ H}\cdot\text{m}^{-1}$, $\mu_B = 9.274 \times 10^{-24}\text{ J}\cdot\text{T}^{-1}$, $c = 2.998 \times 10^8\text{ m}\cdot\text{s}^{-1}$, $1\text{ eV} = 1.602 \times 10^{-19}\text{ J}$, $1\text{ amu} = 1.661 \times 10^{-27}\text{ kg}$.
\vspace{2pt}
\rule{\textwidth}{0.5pt}

\vspace{4pt}
% =============================================================
% PART A: PRE-MIDTERM SYLLABUS (15 MARKS - 30%)
% =============================================================
\begin{center}
    {\large\bfseries PART A --- PRE-MIDTERM SYLLABUS [15 MARKS]}
\end{center}
\vspace{-4pt}

\noindent\textbf{Q1. [Crystal Structure, Reciprocal Lattice \& X-ray Diffraction]} \hfill \textbf{[CO1]}
\begin{enumerate}[label=(\alph*),leftmargin=1.8em,itemsep=3pt]
    \item \textbf{Structure Factor and Systematic Extinctions:}
    \begin{enumerate}[label=(\roman*),itemsep=1.5pt]
        \item Define the geometric structure factor $F_{hkl}$ for a unit cell containing basis atoms at positions $\mathbf{r}_j$.
        \item For a Face-Centered Cubic (FCC) lattice with identical atoms having atomic scattering factor $f$, evaluate $F_{hkl}$ and deduce the systematic reflection selection rules. Prove that reflections with mixed Miller indices (some even, some odd) vanish identically. \hfill\textbf{[1, 3 Marks]}
    \end{enumerate}

    \item \textbf{Powder X-ray Diffraction Analysis:}
    A powder diffraction pattern of polycrystalline Gold (FCC structure, lattice parameter $a = 4.08\text{ \AA}$) is recorded using monochromatic Cu~$K_\alpha$ radiation ($\lambda = 1.5406\text{ \AA}$).
    \begin{enumerate}[label=(\roman*),itemsep=1.5pt]
        \item Determine the Miller indices $(hkl)$ of the first two diffraction peaks observed at the lowest Bragg angles.
        \item Calculate their corresponding interplanar spacings $d_{hkl}$ (in \AA) and the respective glancing Bragg angles $\theta$ (in degrees). \hfill\textbf{[1, 2 Marks]}
    \end{enumerate}
\end{enumerate}

\vspace{4pt}
\noindent\textbf{Q2. [Lattice Vibrations, Phonon Transport \& Free Electron Gas]} \hfill \textbf{[CO2, CO3]}
\begin{enumerate}[label=(\alph*),leftmargin=1.8em,itemsep=3pt]
    \item \textbf{Thermal Conductivity and Phonon Scattering:}
    The lattice thermal conductivity in an insulator is expressed as $\kappa_{ph} = \frac{1}{3} C_v v \ell_{ph}$.
    \begin{enumerate}[label=(\roman*),itemsep=1.5pt]
        \item Detail the physical mechanisms governing the temperature variation of $\kappa_{ph}(T)$ in three regimes: $T \to 0\text{ K}$, intermediate temperatures ($T \sim 0.1\,\Theta_D$), and high temperatures ($T \gg \Theta_D$).
        \item Differentiate between Normal ($N$) and Umklapp ($U$) three-phonon scattering processes. With the help of a reciprocal space vector diagram, explain why Umklapp processes are indispensable for establishing a non-zero thermal resistance in pure dielectric crystals. \hfill\textbf{[3, 2 Marks]}
    \end{enumerate}

    \item \textbf{Wiedemann-Franz Relation and Electronic Thermal Conduction:}
    \begin{enumerate}[label=(\roman*),itemsep=1.5pt]
        \item State the Wiedemann-Franz law. Using the Sommerfeld quantum theory of metals, write the analytical expression for the Lorenz number $L = \frac{\kappa_e}{\sigma T}$ and compute its theoretical numerical value.
        \item Copper has an electrical conductivity $\sigma = 5.90 \times 10^7\text{ }\Omega^{-1}\cdot\text{m}^{-1}$ at $T = 300\text{ K}$. Estimate its electronic thermal conductivity $\kappa_e$ at this temperature. \hfill\textbf{[1, 2 Marks]}
    \end{enumerate}
\end{enumerate}

% =============================================================
% PART B: POST-MIDTERM SYLLABUS (35 MARKS - 70%)
% =============================================================
\vspace{6pt}
\begin{center}
    {\large\bfseries PART B --- POST-MIDTERM SYLLABUS [35 MARKS]}
\end{center}
\vspace{-4pt}

\noindent\textbf{Q3. [Bloch's Theorem, Nearly Free Electron Model \& Tight-Binding Theory]} \hfill \textbf{[CO4]}
\begin{enumerate}[label=(\alph*),leftmargin=1.8em,itemsep=4pt]
    \item \textbf{Nearly Free Electron (NFE) Model and Band Gap Formation:}
    \begin{enumerate}[label=(\roman*),itemsep=1.5pt]
        \item State Bloch's theorem for an electron moving in a 1D crystal potential $V(x) = V(x+a)$.
        \item Near the first Brillouin zone boundary ($k \approx \pm G/2 = \pm \pi/a$), the electron wavefunction is represented by a two-component plane-wave expansion:
        \[
        \psi(x) = C_1 e^{i k x} + C_2 e^{i(k - G)x}, \quad \text{with } G = \frac{2\pi}{a}.
        \]
        For a periodic potential $V(x) = 2V_1 \cos(Gx)$, set up the $2 \times 2$ secular determinant from the time-independent Schr\"odinger equation.
        \item Solve this secular equation at the exact zone boundary $k = G/2$ and prove that an energy band gap of width $E_g = 2|V_1|$ opens between the two standing-wave states.
        \item Sketch the resulting dispersion curve $E(k)$ across the first Brillouin zone in the \textit{reduced zone scheme}, clearly labeling the band gap and energy extrema. \hfill\textbf{[1, 2, 2, 1 Marks]}
    \end{enumerate}

    \item \textbf{Tight-Binding Approximation (LCAO):}
    \begin{enumerate}[label=(\roman*),itemsep=1.5pt]
        \item For a one-dimensional linear atomic chain with lattice constant $a$, consider an electron occupying tightly bound atomic $s$-orbitals. Assuming only nearest-neighbor overlap integral $t$ and on-site energy $\epsilon_0$, derive the energy dispersion relation:
        \[
        E(k) = \epsilon_0 - 2t \cos(ka).
        \]
        \item Find the total bandwidth $W$ of this energy band.
        \item Find the effective mass $m^*$ at the band bottom ($k = 0$) and top ($k = \pi/a$). \hfill\textbf{[2, 1, 1 Marks]}
    \end{enumerate}
\end{enumerate}

\vspace{4pt}
\noindent\textbf{Q4. [Electron Dynamics in Periodic Lattices \& Semiconductor Physics]} \hfill \textbf{[CO4]}
\begin{enumerate}[label=(\alph*),leftmargin=1.8em,itemsep=4pt]
    \item \textbf{Band Electron Dynamics, Negative Effective Mass, and Holes:}
    \begin{enumerate}[label=(\roman*),itemsep=1.5pt]
        \item From the definition of group velocity $v_g = \frac{1}{\hbar}\frac{dE}{dk}$ and the acceleration caused by an applied electric field $\mathcal{E}$, derive the fundamental effective mass equation:
        \[
        m^* = \hbar^2 \left( \frac{d^2E}{dk^2} \right)^{-1}.
        \]
        \item Explain the physical significance of a negative effective mass near the upper edge of an energy band. Explain how an unfilled electron state in an otherwise full valence band is treated as a quasiparticle ``hole'' with charge $+e$ and positive mass $+|m^*|$.
        \item Define \textit{Bloch oscillations}. Why can they not be observed in typical bulk metals under ambient conditions? \hfill\textbf{[1, 2, 1 Marks]}
    \end{enumerate}

    \item \textbf{Semiconductor Statistics and Carrier Concentrations:}
    \begin{enumerate}[label=(\roman*),itemsep=1.5pt]
        \item In a 3D intrinsic semiconductor with band gap $E_g = E_c - E_v$, derive the temperature-dependent intrinsic carrier concentration:
        \[
        n_i(T) = \sqrt{N_c N_v}\, \exp\left(-\frac{E_g}{2 k_B T}\right),
        \]
        where $N_c = 2\left(\frac{2\pi m_e^* k_B T}{h^2}\right)^{3/2}$ and $N_v = 2\left(\frac{2\pi m_h^* k_B T}{h^2}\right)^{3/2}$.
        \item Derive the expression for the intrinsic Fermi level $E_{Fi}$ as a function of temperature:
        \[
        E_{Fi} = \frac{E_c + E_v}{2} + \frac{3}{4} k_B T \ln\left( \frac{m_h^*}{m_e^*} \right).
        \]
        \item Pure Silicon at $300\text{ K}$ has $E_g = 1.12\text{ eV}$, $m_e^* = 1.08\,m_e$, and $m_h^* = 0.56\,m_e$. Calculate $N_c$, $N_v$, and $n_i$ (in $\text{m}^{-3}$ and $\text{cm}^{-3}$). Find the energy shift of $E_{Fi}$ from midgap. \hfill\textbf{[1, 1, 2 Marks]}
    \end{enumerate}
\end{enumerate}

\vspace{4pt}
\noindent\textbf{Q5. [Dielectric Properties, Local Field \& Optical Responses of Solids]} \hfill \textbf{[CO5]}
\begin{enumerate}[label=(\alph*),leftmargin=1.8em,itemsep=4pt]
    \item \textbf{Local Electric Field and Clausius-Mossotti Relation:}
    \begin{enumerate}[label=(\roman*),itemsep=1.5pt]
        \item Differentiate between the macroscopic average electric field $\mathbf{E}$ and the microscopic local field $\mathbf{E}_{loc}$ acting at a lattice site in a dielectric.
        \item By evaluating the polarization charge on a fictitious Lorentz spherical cavity in a cubic dielectric, show that:
        \[
        \mathbf{E}_{loc} = \mathbf{E} + \frac{\mathbf{P}}{3\epsilon_0}.
        \]
        \item Combining the microscopic dipole response $\mathbf{P} = N \alpha \mathbf{E}_{loc}$ with the macroscopic constitutive relation $\mathbf{D} = \epsilon_0 \epsilon_r \mathbf{E} = \epsilon_0 \mathbf{E} + \mathbf{P}$, derive the \textit{Clausius-Mossotti equation}:
        \[
        \frac{\epsilon_r - 1}{\epsilon_r + 2} = \frac{N \alpha}{3\epsilon_0}.
        \]
        \item Explain what is meant by the ``polarization catastrophe'' and how this instability triggers the ferroelectric state with spontaneous polarization. \hfill\textbf{[1, 1, 2, 1 Marks]}
    \end{enumerate}

    \item \textbf{Optical Properties and Plasma Oscillations of Conduction Electrons:}
    \begin{enumerate}[label=(\roman*),itemsep=1.5pt]
        \item Neglecting collision damping, set up the equation of motion for a free conduction electron in an oscillating electric field $\mathbf{E}(t) = \mathbf{E}_0 e^{-i\omega t}$. Derive the frequency-dependent dielectric function of the electron gas:
        \[
        \epsilon(\omega) = 1 - \frac{\omega_p^2}{\omega^2}, \quad \text{where } \omega_p = \sqrt{\frac{n e^2}{\epsilon_0 m_e}}.
        \]
        \item Discuss the optical response of metals in terms of the complex refractive index $\tilde{n} = \sqrt{\epsilon(\omega)}$: explain why metals exhibit high reflectivity ($R \approx 1$) for $\omega < \omega_p$ and become transparent to electromagnetic radiation for $\omega > \omega_p$.
        \item An alkali metal has a conduction electron density $n = 6.00 \times 10^{28}\text{ m}^{-3}$. Calculate its plasma frequency $\omega_p$ (in $\text{rad}\cdot\text{s}^{-1}$), cyclic frequency $\nu_p$ (in $\text{Hz}$), and the corresponding ultraviolet cutoff wavelength $\lambda_p$ (in nm). \hfill\textbf{[1, 1, 2 Marks]}
    \end{enumerate}
\end{enumerate}

\vspace{4pt}
\noindent\textbf{Q6. [Exchange Interaction, Ferromagnetism \& Antiferromagnetism]} \hfill \textbf{[CO5]}
\begin{enumerate}[label=(\alph*),leftmargin=1.8em,itemsep=4pt]
    \item \textbf{Weiss Molecular Field Theory of Ferromagnetism:}
    \begin{enumerate}[label=(\roman*),itemsep=1.5pt]
        \item State the physical premise of Weiss's molecular field hypothesis, writing the internal field as $B_{eff} = B_a + \lambda M$, where $\lambda$ is the molecular field parameter.
        \item Using the Curie law $\chi = M / B_{eff} = C / T$ for paramagnetism above the Curie temperature $T_C$, derive the \textit{Curie-Weiss law}:
        \[
        \chi = \frac{C}{T - \theta_C}, \quad \text{with } \theta_C = \lambda C.
        \]
        \item In the absence of an applied magnetic field ($B_a = 0$), show analytically or graphically using the transcendental equation $M/M_s = \tanh\left(\frac{g\mu_B S B_{eff}}{k_B T}\right)$ that a stable spontaneous magnetization $M_s \neq 0$ appears only when $T < T_C$. \hfill\textbf{[1, 1, 2 Marks]}
    \end{enumerate}

    \item \textbf{Heisenberg Exchange Interaction and Antiferromagnetism:}
    \begin{enumerate}[label=(\roman*),itemsep=1.5pt]
        \item Explain why classical magnetic dipole-dipole interactions fail by several orders of magnitude to explain the observed transition temperatures ($T_C \sim 1000\text{ K}$) of ferromagnets.
        \item Explain the quantum origin of the Heisenberg exchange Hamiltonian:
        \[
        \hat{H}_{ex} = -2 \sum_{i < j} J_{ij}\, \mathbf{S}_i \cdot \mathbf{S}_j,
        \]
        clarifying how Coulomb repulsion combined with the Pauli exclusion principle dictates the sign of the exchange integral $J$. Distinguish between ferromagnetic ($J > 0$) and antiferromagnetic ($J < 0$) coupling.
        \item Sketch the temperature dependence of magnetic susceptibility $\chi(T)$ for an antiferromagnet, clearly highlighting the behavior above and below the N\'eel temperature $T_N$ for fields applied parallel ($\chi_\parallel$) and perpendicular ($\chi_\perp$) to the spin axis. \hfill\textbf{[1, 2, 1 Marks]}
    \end{enumerate}
\end{enumerate}

\vspace{5pt}
\hrule
\vspace{3pt}

% =============================================================
% Course Outcome Mapping & Marks Rubric
% =============================================================
\noindent\textbf{Question-wise Course Outcome (CO) Mapping and Bloom's Taxonomy Level:}
\begin{center}
\small
\begin{tabularx}{\textwidth}{@{}l*{8}{>{\centering\arraybackslash}X}@{}}
\toprule
\textbf{Question No.} & \textbf{Q1} & \textbf{Q2} & \textbf{Q3} & \textbf{Q4} & \textbf{Q5} & \textbf{Q6} & \textbf{Total} & \textbf{Wt.} \\
\midrule
\textbf{Syllabus Domain} & Pre-Mid & Pre-Mid & Post-Mid & Post-Mid & Post-Mid & Post-Mid & --- & --- \\
\textbf{Course Outcome} & CO1 & CO2, CO3 & CO4 & CO4 & CO5 & CO5 & CO1--CO5 & --- \\
\textbf{Bloom's Level}$^*$ & Level 2--3 & Level 2--3 & Level 3--4 & Level 3--4 & Level 3--4 & Level 3--4 & --- & --- \\
\textbf{Marks Allocated} & \textbf{7} & \textbf{8} & \textbf{10} & \textbf{8} & \textbf{9} & \textbf{8} & \textbf{50} & \textbf{100\%} \\
\bottomrule
\end{tabularx}
\end{center}

\vspace{-4pt}
\noindent{\footnotesize
$^*$Bloom's Revised Taxonomy: Level 1 = Remembering; Level 2 = Understanding; Level 3 = Applying; Level 4 = Analysing; Level 5 = Evaluating.\\
\textbf{Course Outcomes (PH7302E):} \textbf{CO1:} Differentiate lattice types, reciprocal lattice and crystal diffraction. \textbf{CO2:} Distinguish crystal excitations and explain specific heat. \textbf{CO3:} Apply free electron model and interpret metallic properties. \textbf{CO4:} Explain energy bands and electrical properties of materials. \textbf{CO5:} Explain dielectric, optical, and magnetic properties of materials.
}

\begin{center}
    \vspace{3pt}
    \textbf{*** END OF QUESTION PAPER --- ALL THE BEST ***}
\end{center}

\end{document}
